Con el fin de alcanzar los objetivos descritos anteriormente, resulta necesario delimitar con precisión el alcance del sistema y concretar qué áreas de actividad se pretenden cubrir mediante la aplicación propuesta.

El alcance de este proyecto se centra en la creación de una aplicación web orientada a la digitalización de la actividad comercial de un distribuidor de productos capilares profesionales. La aplicación pretende centralizar en una única plataforma la relación operativa entre el distribuidor, sus comerciales y los salones clientes, sustituyendo progresivamente procesos dispersos basados en hojas de cálculo, mensajería informal, documentación no centralizada y registros manuales.

La solución propuesta contempla una experiencia web accesible desde distintos dispositivos, con una interfaz adaptada a los principales perfiles de usuario y con soporte básico de instalación como aplicación web progresiva. Además, el sistema se apoya en tecnologías web actuales, persistencia relacional y servicios externos de apoyo para cubrir necesidades como almacenamiento de recursos gráficos, geolocalización, cálculo de rutas, representación de etiquetas y comunicaciones.

\subsection{Alcance funcional y organizativo}
Desde el punto de vista funcional y organizativo, el sistema afecta principalmente al distribuidor administrador, a los comerciales y a los clientes profesionales. A continuación, se detallan las áreas principales incluidas en el alcance del proyecto:

\begin{enumerate}
    \item \textbf{Interfaz para el distribuidor administrador:} el administrador dispondrá de una interfaz desde la que podrá gestionar usuarios, solicitudes de alta, clientes profesionales, comerciales, catálogo, pedidos, comunicaciones, auditoría y configuración general de la aplicación.

    \item \textbf{Interfaz para los comerciales:} los comerciales podrán consultar sus clientes asignados, planificar visitas, organizar rutas, acceder al catálogo, preparar pedidos, consultar repartos, registrar pagos y revisar comunicaciones relevantes para su actividad diaria.

    \item \textbf{Interfaz para los clientes profesionales:} los salones clientes podrán acceder al catálogo técnico-comercial, consultar información vinculada a sus pedidos y repartos, mantener su perfil, revisar comunicaciones comerciales y utilizar herramientas de apoyo relacionadas con su actividad profesional.

    \item \textbf{Gestión de clientes, visitas y rutas:} la aplicación permitirá centralizar los datos de los clientes profesionales, asignarlos a comerciales, registrar visitas, organizar rutas comerciales y facilitar el seguimiento de la relación entre el distribuidor y cada salón.

    \item \textbf{Gestión del catálogo técnico-comercial:} el sistema incluirá la administración y consulta de productos, categorías, líneas, subcategorías, recursos asociados, cartas de color y referencias de tonalidad, de forma que el catálogo pueda utilizarse tanto con finalidad comercial como técnica.

    \item \textbf{Gestión de pedidos, repartos y cobros:} la aplicación cubrirá el flujo operativo de pedido, desde la preparación inicial hasta la entrega por comercial o agencia, incluyendo la organización de repartos, la generación de etiquetas y el registro de pagos parciales o totales.

    \item \textbf{Seguimiento técnico del salón:} el sistema permitirá registrar información técnica asociada a servicios realizados en los clientes profesionales, incluyendo fichas de trabajo, productos utilizados, tonalidades, resultados visuales, plantillas reutilizables y borradores de comunicación técnica.

    \item \textbf{Comunicación comercial y fidelización:} la aplicación contemplará promociones, descuentos, formaciones, avisos internos, recordatorios y comunicaciones multicanal, con el objetivo de mejorar la relación entre el distribuidor, sus comerciales y los salones clientes.

    \item \textbf{Administración, seguridad y trazabilidad:} el alcance incluye el control de acceso por roles, la protección de rutas y operaciones sensibles, el registro de eventos relevantes, la auditoría administrativa y la configuración de aspectos transversales del sistema.
\end{enumerate}

\subsection{Alcance técnico y de integración}
Desde el punto de vista técnico, el proyecto se define como una aplicación web \textit{full-stack} basada en \texttt{Next.js} \citep{nextjsdocs2026}, con persistencia en \texttt{PostgreSQL} \citep{postgresdocs2026} y una capa de acceso a datos apoyada en \texttt{TypeORM} \citep{typeormentities2026,typeormmigrations2026}. Esta base permite mantener una única aplicación para los distintos perfiles de usuario, compartir reglas de negocio y conservar la información del sistema de forma estructurada.

El alcance técnico contempla también el uso de servicios externos de apoyo. Entre ellos se incluyen almacenamiento de imágenes y recursos gráficos mediante \textit{Cloudinary} \citep{cloudinaryupload2026}, geocodificación y cálculo de rutas mediante servicios basados en \textit{OpenStreetMap} y \textit{OSRM} \citep{nominatimapi2026,osrmdocs2026}, y generación de códigos QR para etiquetas operativas mediante \textit{QuickChart} \citep{quickchartqr2026}. Estos servicios no sustituyen la lógica principal de negocio, sino que actúan como soporte para funcionalidades concretas de la aplicación.

La solución se plantea además como una aplicación accesible desde ordenadores, tabletas y teléfonos móviles, incorporando una base de aplicación web progresiva mediante manifiesto web, recursos instalables y \textit{service worker}. Este enfoque permite acercar la experiencia de uso a la de una aplicación nativa sin exigir una distribución específica a través de tiendas de aplicaciones \citep{webdev-pwa}.

Por último, el proyecto contempla una línea de integración con herramientas de gestión empresarial utilizadas en el contexto del distribuidor, especialmente Factusol \citep{sdelsol-factusol}, apoyándose en flujos de automatización externos como \texttt{n8n} cuando sea necesario \citep{n8n-webhook,n8n-http-request}. No obstante, la aplicación no pretende sustituir a un ERP ni asumir por completo los procesos oficiales de facturación, contabilidad o gestión fiscal, sino actuar como una capa operativa propia para la relación comercial diaria y dejar preparada su posible conexión con sistemas externos.
